\section{模型的检验与灵敏度分析}

\subsection{模型检验}

表~\ref{tab:checks} 汇总了全文的检验。“互相验证”只用于\emph{同一问题、不同假设}的方法对，判据取排序或数值的一致性；上下游分工衔接的方法之间不存在验证关系，本文也不作此类声称。

\begin{table}[!htbp]
\centering\small
\caption{全文检验清单}
\label{tab:checks}
\setlength{\tabcolsep}{3pt}
\begin{tabular}{cp{4.3cm}p{2.2cm}p{2.8cm}p{4.3cm}}
\toprule
问题 & 检验对 & 类型 & 判据 & 结果 \\
\midrule
\multirow{6}{*}{一}
& CRITIC--Huber vs 加权均值/PCA/TOPSIS & 方法一致性 & 样本级 Spearman & $0.9936/0.8166/0.9588$ \\
& A1 对 A2/A3 对应域全量 & 分布比较 & KS 检验 & arxiv/github：$p=0.313/0.220$ \\
& 六来源分歧诊断 & 冻结规则复检 & 候选比例 & 高分歧：$5.00\%/4.73\%/4.99\%$；多来源对立：$1.36\%/1.58\%/1.77\%$ \\
& clr+Huber / 裸 OLS / GBDT & 1M 留出集拟合 & 平均逐域 $R^2$ & $0.772/0.654/0.897$ \\
& 配比质量交互 & bootstrap 稳健性 & 区间不含零的域数 & $6/13$，系数正负并存 \\
& 1M $\to$ 60M $\to$ 1B & 冻结模型检验 & 排名 Spearman；$R^2$ & $0.523/0.467/0.677$；$0.772/-7.69/-604.66$ \\
\midrule
\multirow{3}{*}{二}
& B6 拟合 / B7 新增点 & 半合成留出检验 & RMSE & $0.0515/0.0416$ \\
& B7 仿射形状（全 450 点） & 形状检查，非独立检验 & $R^2$；二次项 $t$ & $0.9974$；$-0.35$ \\
& B4/B5 对经典基线 & 跨来源趋势对照 & Spearman & $0.983/0.959$；仅检验经典 $N,D$ 趋势 \\
\midrule
\multirow{3}{*}{三}
& 数值解 vs 闭式 $N^*$ & 互相验证 & 相对误差 & $\le1.83\times10^{-7}$（15 组） \\
& 求解器 vs 网格穷举 & 互相验证 & Loss 差、区制 & $\le1.6\times10^{-5}$，区制一致 \\
& $C_{\mathrm{crit}}$ 三口径 & 互相验证 & 极差 & 主口径 $\le0.05\dex$ \\
\midrule
\multirow{4}{*}{四}
& C8 重算 vs C1 & 口径核验 & Pearson & $\ge0.989$（MATH $0.801$） \\
& sigmoid vs 线性桥接 & 形式选择 & 留一 RMSE & $8.08$ vs $8.28$ \\
& 分解 dex/yr vs 窗口回归 & 互相验证 & 同量级 & $0.29$–$0.78$ vs $0.44$–$0.77$ \\
& 预测回测 & 外部检验 & 2025.25 误差 & M1b $-0.78$，线性外推 $+2.65$ \\
\bottomrule
\end{tabular}
\end{table}

不一致的地方同样是结论。问题一的 17 域质量与训练边际价值秩相关为 $-0.211$，只能说明样本内关联较弱，不能据此推断质量与训练价值完全无关；60M/1B 的配比排序相关仍为正，但绝对损失拟合的 $R^2$ 显著为负，不能把排序保持说成绝对预测有效；问题四中结构模型与线性模型相差 32.8pp，本文如实并列，并解释其来源。A12–A15、B6–B8 与 B10 的证据等级（估算/外推、半合成）在全文中始终保留标注，没有升格为实证。

\subsection{灵敏度与稳健性分析}

\textbf{问题一。}六来源候选比例的操作阈值由 A1 每域的 $B(x)$ 第 95 百分位确定，A2/A3 冻结复用。固定高低端立场为两端 20\% 时，阈值改为 P90/P95/P97.5，A1 候选比例为 $10.00\%/5.00\%/2.50\%$，A2 为 $9.19\%/4.73\%/2.00\%$，A3 为 $9.74\%/4.99\%/2.44\%$；筛查比例随操作点变化，不代表真实冲突发生率。当前无人工质量真值，诊断结果未用于更新主 $Q$。等权与 pile\_cc 单域目标分别搜索配比，因目标不同而得到不同候选份额。

\textbf{问题二。}按 B6 拟合的修订质量律参数为 $\kappa_N=0.4297,\kappa_D=0.1395,k_0=0.1417$；B7 新增 90 个半合成点的留出 RMSE 为 $0.0416$。A、B 质量映射采用 $Q_B=Q_A$ 作为条件性主锚点，当前问题一等权目标质量为 $0.6609$；锚点尚未由损失数据识别，质量项区间估计和锚点敏感性仍待补齐。按式~\eqref{eq:sub} 推导的等效替代关系仅是半合成模型内的解析结论。

\textbf{问题三。}对问题二的 500 组 bootstrap 参数，$C_{\mathrm{crit}}$ 的区间宽度在 $\pm0.03\dex$ 以内。$\eta$ 取 $\pm50\%$ 时，2048 档的 $C_{\mathrm{crit}}$ 移动约 $\pm0.012\dex$，131k 档移动约 $\pm0.2\dex$，此时 $L_{\mathrm{ctx}}^{\mathrm{crit}}$ 相应变为 $60000/20000$。$Q_0$ 相图（图~\ref{fig:q3phase}）表明：$Q_0$ 越低，$C_{\mathrm{crit}}$ 越小，内点区间越宽。

\textbf{问题四。}C6 中 Medium 层的权重取 $0.25$–$1.0$ 时，$S(1.85)$ 的变化不到 $0.4$ 分。C3 历史点的权重取 $0.2$ 与 $1.0$ 时，技术进步率分别为 $0.50$ 与 $0.21\dex/\mathrm{yr}$，因此只作敏感性分析。分解换用 2024$\to$2025 与 2024.5$\to$2025.25 两个窗口，M1b 的技术占比为 $7.1\%$ 与 $8.2\%$，M2 为 $30.6\%$ 与 $24.3\%$，“规模为主”的定性结论不变。

\section{模型的评价与推广}

\subsection{模型的优点}

(1) \textbf{证据分层。}半合成与外推数据分别标注其性质和适用边界；问题一的估算表只用于跨尺度敏感性讨论，问题二以 B1 校准、B6 拟合并用 B7 新增点留出检验质量项。

(2) \textbf{口径可追溯。}Q1 的指标方向、归一化、CRITIC 权重和 Huber 阈值均由 A1 拟合并冻结；六来源冲突诊断分别报告来源内与来源间差异，阈值对应明确的筛查预算；配比模型通过留出集比较，结构性转移由 KKT 活跃集定义。

(3) \textbf{解析表达与数值核验。}问题二旧版的四条命题曾与旧求解器数值逐项对照；本轮修订质量律已推导边际效应、替代条件与算力最优式，但对应数值表尚未重算互验。问题三、四的旧接口数值也须待跨问同步后复核。

(4) \textbf{区分评分与诊断。}质量评分给出唯一主分，六来源冲突分析只用于揭示信号分歧和形成复核候选；没有把筛查比例写成真实冲突率，也没有把评分稳健性证据解释成质量准确率。

\subsection{模型的局限与改进方向}

(1) 11 个 inferred 域的质量分依赖桥接函数（五折 $R^2_{\mathrm{cv}}=0.4106$），区间未纳入模型选择及跨域偏差；六个映射域的验证样本数也有限。若能增补这些域的直接质量信号，可进一步检验桥接假设。

(2) 问题一尚无独立人工质量真值；CRITIC--Huber 与其他排序的一致性、极端指标剔除实验和冲突候选复检都不能替代准确性验证。诊断后的人工复核处置与模型更新尚未完成。

(3) 本轮问题一主质量分和领域分已更新；问题二的新质量形式尚未写回代码和接口，问题三、四仍沿用旧质量锚点与标度律。跨问接口及其数值尚未贯通，需重算核对。

(4) 附件 B 的 B6–B8 为半合成数据，修订质量项的 $\kappa_N,\kappa_D,k_0$ 绝对数值不宜直接外推到真实训练体系；B8 与 B6/B7 方向冲突，应用时须说明数据口径和结构结论的适用范围。

(5) 问题三中 $g(Q)$ 的形式与参数取自赛题附录 B，“角点解”与“对数型早两个数量级提质”对成本函数的形式敏感，本文已把三种形式并列给出。

(6) C6 的 High 层只有 7 点且全部位于底噪区，桥接的上升段依赖 Medium 层；问题四的结构口径与线性口径相差 32.8pp，是全文最大的不确定性来源。若能获得更多同验证集的 Loss–Benchmark 对，可以缩小这一分歧。

(7) 预测为情景条件预测，不包含新架构等离散突破。

\subsection{模型的推广}

“基线 + 质量修正”的标度律结构，以及“约束取等后解析消元”的降维方法，可推广到含多种数据处理成本（去重、合成数据、多模态对齐）的资源配置问题。“S 形桥接 + 连续时间项 + Shapley 分解”的框架也适用于其他存在能力饱和的技术演进分析。

\section{结论}

本文围绕“算力约束下如何配置资源以提升大模型能力”，建立了四个递进衔接的模型。

\textbf{问题一}将 22 个原始字段展开为 25 项“越高越好”的指标，以 CRITIC 加权 Huber 中心形成唯一主质量分，并按中位数聚合至领域。六来源诊断显示 A1/A2/A3 的高分歧候选比例约为 5\%，多来源对立候选约为 1.4\%--1.8\%；这些比例是筛查结果，不是人工确认的真实冲突率。A16/A18 将质量桥接至 17 个配比域，clr 回归建立配比与 13 个验证域交叉熵的关系，GBDT 候选配比相对均匀配比的预测降损为 9.915\%。A6--A11 检验表明 1M 模型的配比排序在 60M/1B 上保留部分相关，但绝对损失预测失效；A12--A15 仅用于估算外推边界。17 域质量与训练边际价值的秩相关为 $-0.211$，不支持按质量排序直接决定训练配比。

\textbf{问题二}以 B1 校准经典律，并按修订形式在 B6 拟合、在 B7 新增点留出检验广义质量项 $L=L_0+(1-Q_B)(\kappa_NAN^{-\alpha}+\kappa_DBD^{-\beta}+k_0)$。在 H0--H2 及锚点假设下，A 侧聚合质量与 B 侧质量的映射可取仿射形式；该质量项的证据来自半合成数据，B8 与其方向不一致，且完整配比修正尚未标定。式中可推出边际效应、质量地板、质量—规模有限替代条件和算力最优规模，但尚未同步到问题三、四接口。

\textbf{问题三}把约束优化解析降维为二维无约束问题，以 KKT 活跃集的改变定义结构性转移，并用 $\Phi(C)=1$ 识别临界预算。主口径下，指数型、幂函数型、对数型的 $\log_{10}C_{\mathrm{crit}}$ 为 $20.61/20.48/18.52$，转移为“不提质 $\to$ 提满”的跳变。上下文长度越过 $6/\eta=30000$ 后，注意力成本占主导，最优模型随之缩小。配比与规模决策近似可分离。

\textbf{问题四}用“标度律 + S 形桥接 + 时间技术项”的结构模型分解前沿增长：规模扩张是主导，非规模技术进步占 $13\%$–$23\%$（结构口径；线性口径的 $30\%$–$46\%$ 为上界参考）。在算力增速放缓至 $0.15\dex/\mathrm{yr}$ 的情景下，24 个月的 base 开源前沿中位数为 $44.1$，技术进步在增量中的占比升至 $77\%$。

四问共同说明，\textbf{数据质量与配比是与算力并列的一等资源}：小预算时应把算力用在规模上，跨过临界预算后应把质量提满；算力增长放缓时，能力增长将越来越依赖非规模技术进步。

需要说明，问题三、四的已列数值仍由旧接口计算，未使用本轮修订的质量锚点和广义标度律；这些下游结论须在接口同步、重新计算后再作为新口径结果。

\begin{thebibliography}{99}\small
\bibitem{hoffmann} Hoffmann J, Borgeaud S, Mensch A, et al. Training compute-optimal large language models[C]//Advances in Neural Information Processing Systems 35 (NeurIPS). 2022.
\bibitem{bahri} Bahri Y, Dyer E, Kaplan J, et al. Explaining neural scaling laws[J]. Proceedings of the National Academy of Sciences, 2024, 121(27): e2311878121.
\bibitem{asai} Asai A, et al. Synthesizing scientific literature with retrieval-augmented language models[J]. Nature, 2026, 650: 857--863.
\bibitem{regmix} Xia M, et al. RegMix: Data mixture as regression for language model pre-training[C]//International Conference on Machine Learning (ICML). 2024.
\bibitem{pythia} Biderman S, et al. Pythia: A suite for analyzing large language models across training and scaling[C]//International Conference on Machine Learning (ICML). 2023.
\bibitem{schaeffer} Schaeffer R, Miranda B, Koyejo S. Are emergent abilities of large language models a mirage?[C]//Advances in Neural Information Processing Systems 36 (NeurIPS). 2023.
\bibitem{muennighoff} Muennighoff N, et al. Scaling data-constrained language models[C]//Advances in Neural Information Processing Systems 36 (NeurIPS). 2023.
\bibitem{fineweb} Penedo G, et al. The FineWeb datasets: Decanting the web for the finest text data at scale[C]//NeurIPS Datasets and Benchmarks Track. 2024.
\bibitem{aitchison} Aitchison J. The Statistical Analysis of Compositional Data[M]. London: Chapman \& Hall, 1986.
\bibitem{critic} Diakoulaki D, Mavrotas G, Papayannakis L. Determining objective weights in multiple criteria problems: The CRITIC method[J]. Computers \& Operations Research, 1995, 22(7): 763--770.
\bibitem{huber} Huber P J. Robust estimation of a location parameter[J]. The Annals of Mathematical Statistics, 1964, 35(1): 73--101.
\bibitem{shapley} Shapley L S. A value for $n$-person games[M]//Contributions to the Theory of Games II. Princeton: Princeton University Press, 1953: 307--317.
\end{thebibliography}
